% ===================================================================
%  QUADRO N.º 14 — A Rua. — Os Comerciantes.
%
%  Ceci n'est PAS une transcription : c'est une traduction, faite en
%  2026, en portugais d'aujourd'hui. Voir texto/pt/10-quadro-01.tex
%  pour la consigne, qui vaut pour les seize fichiers.
% ===================================================================

%%K t14-apar-1 apar
\VUsaut{4.0mm}
\VUtitre{15.0pt}{60}{QUADRO N.º 14}
\VUsaut{3.0mm}

%%K t14-tit-1 sub
\VUpk{11.4pt}{40}{A Rua. --- Os Comerciantes.}
\VUsaut{3.0mm}

%%K t14-01-1 p
1. --- O meu amigo João, que mora no campo, veio passar três dias com a
minha família. \VUgras{Até} agora nunca tinha tido ocasião de ver uma
grande cidade, de modo que passamos todos os dias a passear e eu
explico-lhe o que não conhece.

%%K t14-02-1 p
2. --- Primeiro fomos ver o \VUgras{observatório} \textsuperscript{(1)}, que o
interessou muitíssimo. Ao voltar pela \VUgras{rua} \textsuperscript{(3)} onde estão
os \VUgras{banhos turcos} \textsuperscript{(2)}, cruzámo-nos com um
\VUgras{enterro} \textsuperscript{(4)}. À frente do \VUgras{cortejo fúnebre}
caminhava o \VUgras{clero}, diante do \VUgras{carro fúnebre} no qual se
tinha colocado o \VUgras{caixão}. Muitos amigos seguiam a família de
\VUgras{luto}.

%%K t14-02-2 p
Quando o cortejo passou, atravessámos a rua diante da grande
\VUgras{cervejaria} \textsuperscript{(5)}, onde muitos \VUgras{clientes} bebiam
\VUgras{cerveja} na esplanada, ao abrigo da \VUgras{marquise}
\textsuperscript{(6)} envidraçada.

%%K t14-03-1 p
3. --- João ficou intrigado com o \VUgras{brasão de armas} colocado entre
a loja do \VUgras{comerciante de vinhos e licores} \textsuperscript{(7)} e a do
\VUgras{vendedor de música} \textsuperscript{(10)}. Perguntou-me o que significava;
disse-lhe que assinalava o \VUgras{consulado} \textsuperscript{(8)} inglês, e
mostrei-lhe o \VUgras{cônsul} \textsuperscript{(9)} de pé na varanda.

%%K t14-tit-2 sub
\VUpk{10.2pt}{40}{A ver as montras.}
\VUsaut{3.0mm}

%%K t14-04-1 p
4. --- Naturalmente detivemo-nos diante da exposição de
\VUgras{instrumentos de música}, \VUgras{harmónios}, \VUgras{órgãos} e
\VUgras{pianos}; olhámos as capas ilustradas das \VUgras{partituras} e as
\VUgras{peças} para piano, e admirámos a \VUgras{lira} de madeira dourada
que servia de tabuleta.

%%K t14-05-1 p
5. --- As nuvens de pó que levantavam os \VUgras{varredores}
\textsuperscript{(11)} e a \VUgras{varredoura} \textsuperscript{(12)} obrigaram-nos a
passar depressa diante das belas \VUgras{mobílias} do \VUgras{mobilista}
\textsuperscript{(13)} e da exposição de \VUgras{fogões} e \VUgras{candeeiros} do
\VUgras{ferrageiro} \textsuperscript{(14)}.

%%K t14-06-1 p
6. --- Naquele sítio, com efeito, tivemos de deixar o passeio, pois estava
obstruído pelos volumes que se descarregavam de um \VUgras{carro de
mudanças} \textsuperscript{(16)} para os levar a uma \VUgras{loja}
\textsuperscript{(15)} que acabava de ser arrendada.

%%K t14-06-2 p
Isso impediu-nos de ver as belas carruagens do \VUgras{carroceiro}
\textsuperscript{(17)}, mas não que parássemos a ver o \VUgras{embalador}
\textsuperscript{(19)} que fazia uma caixa para o seu vizinho, o \VUgras{vendedor
de bicicletas e automóveis} \textsuperscript{(18)}. Depois, como já era bastante
tarde, voltámos a casa.

%%K t14-tit-3 sub
\VUpk{10.2pt}{40}{Um passeio pela cidade.}
\VUsaut{3.0mm}

%%K t14-07-1 p
7. --- No dia seguinte a minha mãe propôs que João tirasse um retrato, e
pediu-me que o acompanhasse a casa do \VUgras{fotógrafo} \textsuperscript{(20)}.
Depois do almoço fomos ao \VUgras{estúdio} do artista, onde tivemos de
esperar um bom bocado. Para passar o tempo olhámos os \VUgras{retratos}
\textsuperscript{(22)} pendurados na parede.

%%K t14-07-2 p
Quando chegou a nossa vez a coisa não durou muito, e assim que o meu
amigo acabou de posar diante da \VUgras{máquina} \textsuperscript{(21)}, descemos.

%%K t14-08-1 p
8. --- No primeiro andar vimos, pela porta aberta do \VUgras{cabeleireiro}
\textsuperscript{(23)}, o oficial que cortava o cabelo a um cliente.

%%K t14-08-2 p
Já de novo na rua, passámos diante da \VUgras{tabacaria} \textsuperscript{(24)}. A
João divertiram-no tanto a \VUgras{lanterna} \textsuperscript{(25)} vermelha e o
\VUgras{grande cachimbo} que servia de \VUgras{tabuleta} \textsuperscript{(26)} que
esbarrou com dois velhos que conversavam a tomar uma \VUgras{pitada de
rapé} \textsuperscript{(27)}.

%%K t14-tit-4 sub
\VUpk{10.2pt}{40}{A pastelaria.}
\VUsaut{3.0mm}

%%K t14-09-1 p
9. --- As \VUgras{notas} e as \VUgras{moedas} de todos os países expostas na
\VUgras{montra} do \VUgras{cambista} \textsuperscript{(29)} retiveram-nos um
momento, mas como era a hora do lanche entrámos na \VUgras{pastelaria}
\textsuperscript{(31)} sem nos determos mais.

%%K t14-10-1 p
10. --- Diante de todas as \VUgras{guloseimas} que havia na loja ---
\VUgras{bolos, pão de especiarias, bolachas} e \VUgras{doces} de todas as
espécies --- mal conseguíamos escolher. Por fim João decidiu-se pelos
\VUgras{bolos de creme}, e eu tomei só uma pequena \VUgras{tarte de
cerejas}, pois os doces \VUgras{dão-me dor de dentes} e não tenho nenhuma
vontade de visitar demasiadas vezes o \VUgras{dentista} \textsuperscript{(30)}.

%%K t14-tit-5 sub
\VUpk{10.2pt}{40}{A escrita à máquina.}
\VUsaut{3.0mm}

%%K t14-11-1 p
11. --- Para mostrar ao meu amigo uma \VUgras{máquina de escrever}, de que
tinha ouvido falar mas que nunca tinha visto, entrámos no
\VUgras{escritório} \textsuperscript{(32)} do meu \VUgras{primo} \textsuperscript{(34)}.
Encontrámo-lo a ditar as suas cartas à sua \VUgras{secretária}
\textsuperscript{(35)}, que é \VUgras{estenógrafa}. Mal terminada uma carta, um
\VUgras{dactilógrafo} \textsuperscript{(33)} copiava-a à máquina. Como não
queríamos interromper o trabalho do meu parente, ficámos com ele só uns
minutos.

%%K t14-12-1 p
12. --- Por baixo do escritório do meu primo vive um grande
\VUgras{relojoeiro e joalheiro} \textsuperscript{(37)} que é além disso ourives. A
sua esplêndida loja contém muitos \VUgras{relógios de parede, relógios de
mesa, relógios de bolso}, \VUgras{correntes} e \VUgras{joias} valiosas, como
\VUgras{colares de pérolas}, \VUgras{pulseiras} de ouro, \VUgras{brincos} e
\VUgras{anéis} com \VUgras{pedras preciosas} e \VUgras{diamantes}. Uma das
montras é dedicada por inteiro à \VUgras{ourivesaria} e contém uma grande
coleção de \VUgras{talheres} de prata.

%%K t14-13-1 p
13. --- Ao dobrar a esquina da rua, reparei que tinham mudado a
\VUgras{placa} \textsuperscript{(36)} que está junto ao \VUgras{número} da casa. Ia
dizê-lo em voz alta quando se ouviu o alegre som de uma banda militar.
Deitámos a correr para o regimento, sem pensar que ia passar diante das
lojas do \VUgras{chapeleiro} \textsuperscript{(39)} e do \VUgras{luveiro}
\textsuperscript{(40)}, e que teríamos estado igualmente bem situados para o ver
desfilar ficando diante da montra do \VUgras{oculista} \textsuperscript{(38)}, cheia
de \VUgras{pincenês, óculos} e \VUgras{binóculos}.

%%K t14-tit-6 sub
\VUpk{10.2pt}{40}{O desfile.}
\VUsaut{3.0mm}

%%K t14-14-1 p
14. --- À frente do \VUgras{regimento} \textsuperscript{(41)} marchava o
\VUgras{tambor-mor} \textsuperscript{(42)}, seguido dos \VUgras{tambores}
\textsuperscript{(43)}, dos \VUgras{corneteiros} \textsuperscript{(44)} e de toda a
\VUgras{banda}. O \VUgras{general} \textsuperscript{(45)}, que estava na praça com o
seu ajudante de campo, tinha parado junto ao \VUgras{repuxo}
\textsuperscript{(49)} do \VUgras{chafariz} \textsuperscript{(48)} para ver o
\VUgras{desfile}.

%%K t14-14-2 p
\VUgras{Os oficiais do estado-maior} \textsuperscript{(46)}, o \VUgras{coronel} e o
\VUgras{tenente-coronel} saudavam-no com a espada. \VUgras{Os majores} iam à
frente dos seus \VUgras{batalhões} \textsuperscript{(47)} e cada \VUgras{capitão}
comandava a sua \VUgras{companhia}, enquanto os \VUgras{tenentes}, os
\VUgras{alferes} e os \VUgras{sargentos} ocupavam os seus postos habituais
junto aos seus homens.

%%K t14-15-1 p
15. --- Muitos transeuntes se tinham juntado no \VUgras{cruzamento}
\textsuperscript{(50)}; os comerciantes --- o \VUgras{guarda-chuveiro}
\textsuperscript{(51)}, o \VUgras{tintureiro} \textsuperscript{(53)}, o \VUgras{armeiro}
\textsuperscript{(54)} --- saíram para ver os \VUgras{soldados}. Todos os veículos
--- \VUgras{carruagens de praça} \textsuperscript{(52)}, \VUgras{carruagens
particulares} \textsuperscript{(57)} e também o \VUgras{carro de rega}
\textsuperscript{(56)} --- se encostaram ao passeio.

%%K t14-tit-7 sub
\VUpk{10.2pt}{40}{Cenas da rua.}
\VUsaut{3.0mm}

%%K t14-15-2 p
Um \VUgras{caixeiro-viajante} \textsuperscript{(55)} que levava na mão as suas
\VUgras{malas de amostras} atravessou a rua a bom passo, e as pessoas
sentadas na \VUgras{imperial} \textsuperscript{(59)} do \VUgras{ómnibus}
\textsuperscript{(58)} voltaram-se para gozar o espetáculo. Um pobre
\VUgras{cantor de rua} \textsuperscript{(60)} com o seu \VUgras{realejo}
\textsuperscript{(61)} teve de interromper a sua \VUgras{canção}, pois o ruído dos
tambores abafava a sua voz.

%%K t14-16-1 p
16. --- Quando o regimento chegou diante da \VUgras{loja de tecidos}
\textsuperscript{(64)}, as \VUgras{costureiras} \textsuperscript{(63)} e os
\VUgras{alfaiates} \textsuperscript{(62)} deixaram as suas \VUgras{máquinas de
costura} e a sua tarefa para se debruçarem à janela. \VUgras{O comerciante}
\textsuperscript{(65)}, que acabava de pôr \VUgras{fatos} \textsuperscript{(66)} novos na
sua \VUgras{montra}, teve de retirar as peças de \VUgras{pano}
\textsuperscript{(67)} e de \VUgras{linho} expostas no passeio, junto à \VUgras{boca
de esgoto} \textsuperscript{(68)}, com medo que as pessoas as deitassem abaixo;
até um velho \VUgras{bufarinheiro} \textsuperscript{(69)}, a quem uma porteira
regateava umas meias, teve de se meter num portal para acabar a sua
venda.

%%K t14-16-2 p
Acompanhámos o regimento até ao quartel \VUgras{e}, muito contentes com o
nosso dia, voltámos a casa à hora do jantar.

%%K t14-tit-8 sub
\VUpk{10.2pt}{40}{Ofícios e profissões diversos.}
\VUsaut{3.0mm}

%%K t14-17-1 p
17. --- No dia seguinte o meu pai ofereceu-se para nos levar a visitar a
tipografia do jornal local. Aceitámos com entusiasmo.

%%K t14-17-2 p
O \VUgras{impressor} é também \VUgras{livreiro} \textsuperscript{(70)},
\VUgras{editor}, \VUgras{papeleiro} e \VUgras{encadernador}. Instalou no
primeiro andar a \VUgras{redação} \textsuperscript{(71)} do jornal, e também a
\VUgras{oficina de gravura}, onde trabalham os \VUgras{litógrafos}
\textsuperscript{(72)} e os \VUgras{gravadores em cobre} \textsuperscript{(73)}, e por fim
a sala de composição, onde estão os \VUgras{tipógrafos} \textsuperscript{(74)}. A
\VUgras{sala das máquinas} \textsuperscript{(75)} propriamente dita, as
\VUgras{prensas} e as diversas máquinas, estão no rés do chão.

%%K t14-tit-9 sub
\VUpk{10.2pt}{40}{João faz muitíssimas perguntas.}
\VUsaut{3.0mm}

%%K t14-18-1 p
18. --- Ao sair da tipografia, João parou, muito intrigado, diante de uma
caixinha de vidro montada sobre um poste em frente da \VUgras{retrosaria}
\textsuperscript{(77)}, e perguntou para que servia. O meu pai explicou-lhe que
era um \VUgras{avisador de incêndios} \textsuperscript{(76)}. Com efeito, tudo na
cidade era uma maravilha para o meu amigo, desde o simples
\VUgras{bebedouro} \textsuperscript{(78)} e o ligeiro \VUgras{triciclo de entregas}
\textsuperscript{(80)} que montava um \VUgras{paquete} \textsuperscript{(79)} de libré,
até ao \VUgras{furgão do correio} \textsuperscript{(81)}.

%%K t14-19-1 p
19. --- Enquanto o meu pai nos deixava um momento para falar com o
\VUgras{recebedor das contribuições} \textsuperscript{(83)}, que se tinha detido a
conversar com um \VUgras{advogado} \textsuperscript{(82)} e um \VUgras{notário}
\textsuperscript{(84)}, entretivemo-nos a seguir o movimento da rua. Passava
diante de nós a gente mais variada: um \VUgras{aprendiz de pasteleiro}
\textsuperscript{(85)}, que levava num cesto um esplêndido \VUgras{bolo}, vinha
atrás de um \VUgras{adelo} \textsuperscript{(86)}; um \VUgras{acendedor de
candeeiros} \textsuperscript{(87)} falava com um \VUgras{gasista}
\textsuperscript{(88)}; uma \VUgras{freira} \textsuperscript{(89)} que levava pela mão uma
orfãzinha voltava ao seu convento.

%%K t14-tit-10 sub
\VUpk{10.2pt}{40}{De volta a casa.}
\VUsaut{3.0mm}

%%K t14-20-1 p
20. --- Justamente quando o meu pai voltava para junto de nós, João soltou
uma gargalhada ao ver as caras tisnadas e a estranha figura de dois
\VUgras{limpa-chaminés} \textsuperscript{(91)} negros de fuligem.

%%K t14-21-1 p
21. --- Eu queria comprar uma planta para a minha mãe à \VUgras{florista}
\textsuperscript{(92)}, mas o meu pai observou que pesaria muito para a levar e
que seria melhor comprar umas \VUgras{laranjas}. Aproximámo-nos da
\VUgras{vendedeira da banca} \textsuperscript{(94)}, e quando a \VUgras{aia}
\textsuperscript{(93)}, que estava a escolher algumas para o seu menino, acabou a
sua compra, serviram-nos.

%%K t14-22-1 p
22. --- De volta a casa encontrámos vários conhecidos: uma velha amiga da
minha família, apoiada no braço da sua \VUgras{dama de companhia}
\textsuperscript{(95)}, o \VUgras{engenheiro} \textsuperscript{(96)} de pontes e estradas,
e uma das minhas primas com a sua \VUgras{ama-seca} \textsuperscript{(98)}.

%%K t14-22-2 p
Depois passámos junto à \VUgras{modista} \textsuperscript{(97)} da minha mãe e a
dois \VUgras{frades} \textsuperscript{(99)} forasteiros, que se aproximaram para
perguntar ao meu pai o caminho da estação.
\VUsaut{6.0mm}
\VUfilet{22mm}
